\addcontentsline{toc}{section}{\textbf{CAPÍTULO VI}}
\addcontentsline{toc}{section}{\textbf{CONCLUSIONES Y RECOMENDACIONES}}
%====================================================
\subsection{Conclusiones}
%====================================================

\begin{itemize}

\item \textbf{Primera:} Se demostró que la aplicación del pipeline de digitalización basado en la arquitectura Transformer TrOCR con \textit{fine-tuning} optimiza significativamente la captura de conocimiento técnico, logrando reducir la tasa de error de caracteres (CER) en fórmulas matemáticas y notación científica compleja de un 34.20\% (utilizando métodos OCR tradicionales) a un 4.15\% estable. Este impacto se traduce directamente en una ganancia neta del 30.05\% en la precisión de renderizado en formato \LaTeX{}, eliminando las barreras de distorsión algebraica y asegurando que los estudiantes de la UNAMBA dispongan de un repositorio digital íntegro, fiel y de alta disponibilidad para el repaso de asignaturas de ciencias exactas.

\item \textbf{Segunda:} Se comprobó que la integración de un middleware basado en la arquitectura de Generación Aumentada por Recuperación (RAG) ejerce una influencia crítica en la calidad del aprendizaje asimilado, al mitigar las alucinaciones factuales de los modelos de lenguaje a un índice asintótico del 1.20\% (en comparación con el 28.50\% presente en asistentes de IA genéricos). El condicionamiento del espacio de inferencia de los LLMs al corpus académico validado del estudiante garantiza la exactitud y veracidad factual de la sumarización abstractiva, lo que impactó de manera directa en el rendimiento de los alumnos al elevar sus calificaciones promedio en exámenes parciales de un rango regular (11.80) a un nivel sobresaliente (15.90).

\item \textbf{Tercera:} Se validó que la estructuración y modelamiento de la información a través de Grafos de Conocimiento, en sinergia con la generación automatizada de evaluaciones formativas (\textit{quizzes}), optimiza significativamente los procesos cognitivos y el aprendizaje colaborativo. Metodológicamente, se constató un antes y un después disruptivo: la automatización de cuestionarios dinámicos de repaso espaciado redujo la sobrecarga cognitiva autopercibida por los estudiantes en 46.30 puntos (según la escala NASA-TLX) y fijó una retención conceptual a largo plazo superior al 93.50\%; asimismo, la topología relacional enlazada en red catalizó la inteligencia colectiva, transformando notas individuales aisladas en un ecosistema compartido que incrementó la transferencia eficiente de material de estudio entre pares de 2.10 a 18.40 ítems semanales.

\end{itemize}

%====================================================
\subsection{Recomendaciones}
%====================================================

\begin{itemize}

\item \textbf{Primera:} Se recomienda a la Escuela Académico Profesional de Ingeniería de Informática y Sistemas de la UNAMBA promover y estandarizar el uso de pipelines basados en visión artificial autorregresiva (como TrOCR) para la preservación digital del patrimonio manuscrito de las asignaturas de ingeniería, expandiendo el dataset local con muestras de caligrafía regional y variables de ruido ambiental para robustecer los procesos de \textit{fine-tuning} ante apuntes de laboratorios físicos o pizarras institucionales.

\item \textbf{Segunda:} Se sugiere implementar mecanismos de actualización y curaduría dinámica sobre los vectores semánticos del middleware RAG, con el fin de anexar de forma automatizada los sílabos vigentes, las guías de práctica y la literatura científica actualizada provista por los docentes. Esto garantizará que el corpus de control relacional no sufra obsolescencia contextual y continúe blindando los resúmenes automáticos contra cualquier sesgo cognitivo o alucinación factual del LLM a lo largo de los semestres académicos.

\item \textbf{Tercera:} Se recomienda migrar y escalar la infraestructura del motor de Grafos de Conocimiento relacional y vectorial hacia un modelo distribuido en la nube, aprovechando el desacoplamiento arquitectónico de SynapTome. Esto permitirá soportar de forma eficiente la concurrencia masiva y síncrona en tiempo real mediante SignalR cuando el ecosistema se democratice e interconecte de manera transversal a otras facultades universitarias, maximizando así el descubrimiento semántico interdisciplinario y potenciando el entorno de aprendizaje compartido entre múltiples comunidades estudiantiles.

\end{itemize}